\chapter*{Conclusion}

This thesis has examined the theoretical foundations and practical applications of vector databases, with a particular focus on approximate nearest neighbor (ANN) search in high-dimensional spaces. Through detailed analysis of ANN paradigms, including hash-based methods (LSH), tree-based methods (Annoy), and graph-based methods (HNSW), the study clarified their operational principles, advantages, and limitations, providing insights that complement existing literature.

The practical investigation highlighted contemporary vector management systems such as FAISS, Milvus, and Weaviate. By generating embeddings from a representative dataset, constructing indices, and executing similarity searches, the thesis demonstrated how these systems can be effectively leveraged in real-world scenarios. This hands-on component validated theoretical concepts and provided a clear workflow for implementing similarity search in high-dimensional spaces.

\section*{Achieved Objectives of the Thesis}

The primary objectives set at the outset have been successfully met:

\begin{itemize}
    \item Established a comprehensive theoretical foundation, covering vector representations, similarity measures, and nearest neighbor search principles.
    \item Analyzed approximate nearest neighbor algorithms, including hash-based (LSH), tree-based (Annoy), and graph-based (HNSW) methods, highlighting their characteristics, advantages, and limitations.
    \item Demonstrated the practical application of vector databases through FAISS, Milvus, and Weaviate, including embeddings storage, index construction, and similarity-based querying, with experiments that bridged theory and practice on real datasets.
\end{itemize}

\section*{Future Directions}

Building on the work presented in this thesis, future research could focus on enhancing performance and scalability through techniques such as multithreading, CPU-level optimizations, and low-level programming, which were beyond the scope of this study. Other promising directions include the development of hybrid ANN algorithms, large-scale distributed vector database architectures, and tighter integration with AI pipelines to improve retrieval efficiency and versatility.